\chapter{Introdução}
\label{cap1_introducao}

O objetivo deste documento é apresentar um modelo básico para a elaboração de Trabalho de Conclusão de Curso (TCC) do Instituto Federal do Norte de Minas Gerais (IFNMG), utilizando \LaTeX{} e a classe \textbf{abnTeX2}, em conformidade com as normas da Associação Brasileira de Normas Técnicas (ABNT), em especial a \texttt{ABNT NBR 14724:2011} (trabalhos acadêmicos) e a \texttt{ABNT NBR 10520:2023} (citações).

A introdução visa oferecer uma visão clara do que será tratado no trabalho e normalmente contempla os seguintes elementos:

\begin{alineas}
	\item \textbf{Contextualização do tema}: apresente o cenário atual do tema abordado, explicando sua relevância e o contexto no qual ele está inserido;
	\item \textbf{Problema de pesquisa}: defina de forma clara qual é o problema que o trabalho pretende abordar ou resolver, seja uma questão ou uma necessidade identificada na área de estudo;
	\item \textbf{Objetivo geral}: descreva o que o trabalho pretende alcançar de forma ampla;
	\item \textbf{Objetivos específicos}: detalhe as etapas ou os aspectos que devem ser abordados para atingir o objetivo geral;
	\item \textbf{Justificativa}: explique a relevância do trabalho, as possíveis contribuições para a academia, o mercado ou a sociedade e sua relação com outros estudos;
	\item \textbf{Metodologia (breve)}: mencione resumidamente os métodos e as abordagens utilizados (por exemplo, estudo de caso, pesquisa bibliográfica, desenvolvimento de um sistema);
	\item \textbf{Estrutura do trabalho}: resuma o conteúdo dos capítulos subsequentes.
\end{alineas}

O restante do documento está organizado da seguinte forma: o Capítulo~\ref{cap2_revisao} apresenta os elementos principais, como referências, figuras, quadros, tabelas, equações e algoritmos; o Capítulo~\ref{cap3_metodologia} descreve a metodologia; o Capítulo~\ref{cap4_resultados} apresenta os resultados; e o Capítulo~\ref{cap5_conclusao} traz as considerações finais. Por último, veja um exemplo de apêndice (Apêndice~\ref{apendiceA}) e de anexo (Anexo~\ref{anexoA}).
